\chapter{Hello, parentheses}
\label{ch:hello}

Scheme has a reputation for parentheses. Here is the secret that makes them
friendly: \textbf{there is only one rule}. A list in parentheses means
``apply the first thing to the rest''. That's it. \code{(+ 1 2)} applies
\code{+} to \code{1} and \code{2}. \code{(fs-read "plants.txt")} applies
\code{fs-read} to a file name. \code{(define x 5)} is a special case you will
meet in a minute, and it looks exactly the same.

\section{Where to type}

You have three places to try Scheme, from lightest to fullest:

\begin{sh}
nodwyn eval -e '(* 6 7)'       # one expression, prints its value
nodwyn repl                    # an interactive loop over the language alone
nodwyn                         # a session; type :eval EXPR at the prompt
\end{sh}

\code{eval} and \code{repl} are the language on its own: no files, no
programs, no model, just the builtins. A session's \code{:eval} adds the
\emph{host primitives} (\code{fs-read}, \code{shell-exec} and friends), and
the guard. Part~II needs only the first two.

\begin{tryit}
\begin{txt}
$ nodwyn repl
nodwyn> (define (sq x) (* x x))
nodwyn> (map sq (iota 5))
(0 1 4 9 16)
\end{txt}
You just defined a function and applied it to the list \code{(0 1 2 3 4)}.
We will take that apart slowly.
\end{tryit}

\section{Numbers}

\begin{scm}
(* 6 7)            ;=> 42
(+ 1 2 3 4)        ;=> 10       most operators take any number of arguments
(- 10)             ;=> -10
(/ 9 3)            ;=> 3        exact when the answer is a whole number
(/ 1 3)            ;=> 0.3333333333333333   a float otherwise
(quotient 17 5)    ;=> 3
(remainder -7 2)   ;=> -1       takes the sign of the first number
(modulo -7 2)      ;=> 1        takes the sign of the second
(expt 2 10)        ;=> 1024
(sqrt 16)          ;=> 4.0
(max 3 9 4)        ;=> 9
(exact 3.9)        ;=> 3        float to integer, toward zero
\end{scm}

There are two kinds of number: 64-bit integers and 64-bit floats. Integers
wrap around on overflow, exactly as in Go, which you will almost never
notice, and there are no fractions or big numbers.

Because the operator always comes first, there is no precedence to
remember. \code{(+ 1 (* 2 3))} is 7 and nothing else.

\section{Naming things: \texttt{define}}

\begin{scm}
(define garden-size 12)
garden-size                    ;=> 12

(define (area width height)
  (* width height))
(area 3 4)                     ;=> 12
\end{scm}

Names can contain almost anything that is not a parenthesis or a space:
\code{water-need}, \code{plant?}, \code{set-level!}, \code{->string} are
all fine. By convention a name ending in \code{?} answers yes or no, and one
ending in \code{!} changes something.

\section{Functions are values: \texttt{lambda}}

\code{(define (area w h) \ldots)} is shorthand. What it really does is:

\begin{scm}
(define area
  (lambda (width height)
    (* width height)))
\end{scm}

\code{lambda} makes a function; \code{define} gives it a name. You can use
a \code{lambda} without naming it at all:

\begin{scm}
((lambda (x) (* x x)) 9)       ;=> 81
\end{scm}

This matters more than it looks. Functions in Scheme are values like
numbers: you can pass them to other functions, return them, keep them in
lists. The whole of Part~IV rests on this: a strategy wraps the loop's
functions in functions of its own.

\section{Deciding: \texttt{if}, \texttt{cond}, \texttt{and}, \texttt{or}}

\begin{scm}
(if (> 3 2) "yes" "no")        ;=> "yes"

(define (water-need plant)
  (cond ((equal? plant "basil") "daily")
        ((equal? plant "cactus") "monthly")
        (else "weekly")))
(water-need "cactus")          ;=> "monthly"
\end{scm}

\textbf{Only \code{\#f} is false.} Everything else is true, including
\code{0}, the empty string and the empty list:

\begin{scm}
(if 0 "zero is true" "false")               ;=> "zero is true"
(if '() "empty list is true" "false")       ;=> "empty list is true"
\end{scm}

\code{and} and \code{or} stop as soon as they know the answer, and return
the value that decided it, not just \code{\#t}:

\begin{scm}
(and 1 2 3)                    ;=> 3
(or #f #f "found")             ;=> "found"
(and)                          ;=> #t
(or)                           ;=> #f
\end{scm}

That makes \code{or} a handy ``default'': \code{(or (lookup x) "unknown")}.

\section{Local names: \texttt{let} and \texttt{let*}}

\begin{scm}
(let ((a 1) (b 2))
  (+ a b))                     ;=> 3

(let* ((a 1)
       (b (+ a 1)))            ; let* sees the names above it
  (* a b))                     ;=> 2
\end{scm}

\code{let} binds all its names at once, so they cannot see each other;
\code{let*} binds them one after another. When in doubt, use \code{let*}.

\section{Quoting: code as data}

If \code{(+ 1 2)} means ``apply \code{+}'', how do you write the list
itself, three elements, without applying anything? You \textbf{quote} it:

\begin{scm}
(quote (+ 1 2))                ;=> (+ 1 2)
'(+ 1 2)                       ;=> (+ 1 2)     the apostrophe is shorthand
'basil                         ;=> basil       a symbol
\end{scm}

A \textbf{symbol} is a name used as a value. You will see symbols as keys
all over \nodwyn: a tool call is a list like \code{((id . "c1") (name .
"fs-read") (args . "\{\ldots\}"))} whose keys \code{id}, \code{name} and
\code{args} are symbols.

\section{When things go wrong}

Errors say who complained and why:

\begin{txt}
(car 5)            nodwyn: car: argument 1: expected pair, got integer 5
nope               nodwyn: eval: unbound variable: nope
((lambda (x) x))   nodwyn: #<procedure>: expected 1 argument(s), got 0
(set! nope 1)      nodwyn: set!: unbound variable: nope
\end{txt}

\code{set!} changes an \emph{existing} name; it will not create one. That
catches typos: \code{(set! gardn-size 13)} is an error, not a new variable.

\section{The types, all of them}

\begin{center}
\begin{tabularx}{\linewidth}{llX}
\toprule
\textbf{type} & \textbf{looks like} & \textbf{notes} \\
\midrule
integer & \code{42} \code{-7} & 64-bit \\
float & \code{2.5} \code{1e3} & always printed with a dot or an exponent \\
boolean & \code{\#t} \code{\#f} & only \code{\#f} is false \\
string & \code{"basil"} & immutable; length counts characters, not bytes \\
symbol & \code{basil} & a name as a value \\
list & \code{(1 2 3)} \code{'()} & chains of pairs; next chapter \\
pair & \code{(a . b)} & two things glued together \\
vector & \code{\#(1 2 3)} & indexed, and the one mutable collection \\
procedure & \code{\#<procedure area>} & functions, yours and builtin \\
handle & \code{\#<cred:class:id>} & a credential; cannot become text \\
void & (prints nothing) & what \code{define} and \code{display} return \\
\bottomrule
\end{tabularx}
\end{center}

Each has a predicate: \code{number?}, \code{integer?}, \code{float?},
\code{boolean?}, \code{string?}, \code{symbol?}, \code{pair?},
\code{list?}, \code{null?}, \code{vector?}, \code{procedure?},
\code{handle?}. There is no separate character type: a one-character
string does that job.

\begin{puzzle}[ 1: temperatures]
Write \code{(c->f c)} that turns Celsius into Fahrenheit, and
\code{(frost? c)} that says whether a temperature is at or below freezing.
\code{(c->f 100)} should be \code{212}, and \code{(c->f 21)} should be
\code{69.8}.\ans{1}
\end{puzzle}

\begin{puzzle}[ 2: seasons]
Write \code{(season month)} for months 1 to 12 that returns
\code{"winter"}, \code{"spring"}, \code{"summer"} or \code{"autumn"} (in
the northern hemisphere, say December to February is winter). Try it with
\code{cond}, then look up \code{case} in Chapter~\ref{ch:forms} and write it
again.\ans{2}
\end{puzzle}
